\begin{table}[htbp]
\centering
\caption{敏感性运行结果汇总（10 组设置相对正式基线；逐对删除只作边际相关诊断，选择类指标不适用）}\label{tab:q2_sensitivity}
\begin{tabular}{p{6.2cm}cccccc}
\toprule
设置 & 化合物层 Jaccard & 模块层 Jaccard & Top-1 改变维度数 & 等级改变维度数 & 第 1 位入选频率 & A/B/C 维度数 \\
\midrule
正式基线（τ=0.6，225 候选，列均值，附件 1，重建图，日期组，x̄ 的秩） & 1.000 & 1.000 & 0 & 0 & 0.599 & 26/22/1 \\
模块阈值 τ=0.5 & 0.980 & 0.980 & 0 & 0 & 0.597 & 26/22/1 \\
模块阈值 τ=0.7 & 1.000 & 1.000 & 0 & 0 & 0.600 & 26/22/1 \\
候选加回 5 个综合指标（230） & 0.871 & 0.871 & 3 & 2 & 0.569 & 24/24/1 \\
候选加回内标（226） & 1.000 & 1.000 & 0 & 0 & 0.599 & 26/22/1 \\
单化合物 NA 最低位并列 & 0.980 & 0.980 & 0 & 0 & 0.598 & 26/22/1 \\
24 款啤酒 80 个冲突单元格改用附件 8；图由替换后的附件 1 逐对重算 & 0.963 & 0.963 & 0 & 1 & 0.605 & 27/21/1 \\
bootstrap 中固定附件 2 图（与主分析同一组日期组索引） & 1.000 & 1.000 & 0 & 0 & 0.599 & 26/22/1 \\
bootstrap 单位改为单款啤酒（250 有放回） & 1.000 & 1.000 & 0 & 1 & 0.635 & 27/21/1 \\
感官目标改用附件 4 的 z 的秩（49 维，按行号与 x̄ 对齐） & 0.531 & 0.569 & 12 & 10 & 0.565 & 21/26/2 \\
逐对删除：只作边际 Spearman 诊断（感官列为 x̄；附件 6 复算取 z） & -- & -- & -- & -- & -- &  \\
\bottomrule
\end{tabular}
\end{table}
